\section{Introduction}
\label{sec:intro}

The modern paradigm of multi-task deep learning serving relies heavily on Parameter-Efficient Fine-Tuning (PEFT) methods, particularly Low-Rank Adaptation (LoRA) \cite{hu2021lora}. By dynamically routing and ensembling specialized adapters on-the-fly, serving systems can process heterogeneous, multi-domain inputs on edge devices without the overhead of hosting multiple monolithic base models \cite{pfeiffer2020adapterhub, wortsman2022model}. However, dynamic ensembling in deep architectures is fundamentally constrained by two core challenges: \emph{routing volatility} (jitter) and \emph{cascading representational drift}.

Stateless dynamic ensembling methods, such as SABLE \cite{sable2025} or SPS-ZCA \cite{spszca2025}, compute ensembling weights independently at each layer using local activation similarities to pre-computed centroids. This formulation lacks temporal context across depth, resulting in high-frequency layer-wise ensembling weight fluctuations or "jitter" \cite{chemmerge2025}. In contrast, stateful continuous-time ensembling methods, like ChemMerge \cite{chemmerge2025}, model ensembling weights as concentrations governed by a system of discretized reaction-decay ordinary differential equations (ODEs). This ODE-based tracking introduces spatial inertia across layers, acting as a natural low-pass filter that smooths trajectories and improves representational consistency \cite{chen2018neural, haber2017stable}.

Despite its advantages, stateful ensembling introduces a severe and fragile heuristic known as \emph{active representation coupling} (feedback warping). Under this scheme, intermediate hidden representations are warped toward active task-specific centroids by a constant feedback step size $\eta$. Under practical serving conditions where task centroids are extracted once at an early layer, setting a constant feedback rate $\eta > 0$ is highly suboptimal. Under steady-state homogeneous streams, constant feedback actually degrades performance monotonically as $\eta$ increases. This occurs because late-stage representations are already highly refined by the expert's residual blocks across depth; warping them toward static early-stage centroids pulls them backward toward the noisier early-stage coordinate space. Conversely, under heterogeneous streams, the kinetics propagation lag means representations are unaligned early on, so constant feedback provides a minor boost to ensembling weight convergence but suffers from the same representational backward-shift in later layers. This dual effect results in a flat or marginally positive response to constant feedback, which fails to capture the optimal performance level of adaptive closed-loop controllers. To avoid these representational mismatches, we require an active closed-loop controller that dynamically regulates the feedback rate on-the-fly.

\begin{figure}[t]
\centering
\includegraphics[width=0.48\textwidth]{coupling_ablation.png}
\caption{The Impact of constant feedback warping ($\eta$) on Joint Mean Accuracy across 10 random seeds. Under steady-state homogeneous streams, any constant feedback rate $\eta > 0$ monotonically degrades performance due to representational backward-shift (pulling highly refined late-layer activations toward static early-stage centroids). Under heterogeneous streams, constant feedback is relatively flat or marginally positive. L-ARC (Ours) completely bypasses these limitations by dynamically adapting the feedback rate on-the-fly via our control-theoretic closed-loop controller, outperforming all constant feedback baselines on both workloads.}
\label{fig:coupling_ablation}
\end{figure}

From a control-theoretic perspective, a constant feedback rate represents an open-loop controller trying to stabilize a highly non-linear, state-dependent dynamical system. In this paper, we argue that empirical heuristics are insufficient, and continuous-depth routing should be stabilized by applying formal control-theoretic principles. To this end, we present \textbf{Lyapunov-Stable Active Representation Coupling (L-ARC)}, a training-free dynamic serving framework centered on \emph{state-shielded continuous kinetics}.

The core engine of L-ARC is \textbf{Entropy-Gated Concentration Gating (ECG-Reset)}, which monitors routing entropy to dynamically freeze continuous ODE kinetics during routing dropouts or failures. To complement this state-space shielding, we introduce a closed-loop \textbf{Lyapunov Feedback Controller} that models representation error as a candidate Lyapunov function. This analytical formulation yields a local \textbf{Dissipation Guard} that computes the local dissipation coefficient $A_b^{(l)}$ on-the-fly, gating off representation warping when the update is non-dissipative ($A_b^{(l)} \le \theta_G = 0.04$) and scaling it adaptively when stable ($A_b^{(l)} > \theta_G$). While active feedback is mathematically elegant, our rigorous empirical profiling reveals a critical trade-off: under clean serving conditions, active feedback is practically redundant ($p = 0.106$) while introducing a relative routing latency overhead. To bypass this, we introduce \textbf{Entropy-Triggered Lyapunov Gating (ET-L-ARC)}, a control-theoretic optimization that evaluates the Dissipation Guard only under moderate routing uncertainty ($0.15 \le H \le 0.95$). This collapses our relative latency overhead under clean serving to just $99.85\%$ (adding only $0.06$ ms of absolute wall-clock latency per sample), making L-ARC exceptionally practical. Under transient failures (Setting C), the dual-shield of ECG-Reset and our Dissipation Guard provides critical, robust shielding---preventing both stateful memory corruption and representation-space semantic collapse. Finally, under Confident Router Bias (Setting D), we propose \textbf{Representation-Agreement State Correction (RASC)}, a dual-loop control mechanism that compares feedforward routing selections with representation-space coordinates, overriding corrupted, biased router rates to completely resolve state-locking failures.

Our key contributions are:
\begin{itemize}
    \item \textbf{Stateful Serving Framework (ECG-Kinetics):} We propose Entropy-Gated Concentration Gating (ECG-Reset) to shield continuous-depth model ensembling, preventing the propagation of memory faults under transient failures.
    \item \textbf{Control-Theoretic Feedback Formulation:} We construct a formal Lyapunov stability framework analyzing representation error across layer depth, presenting active representation feedback as a theoretically elegant boundary case.
    \item \textbf{Analytical Dissipation Guard:} We derive a closed-loop controller that regulates the feedback rate on-the-fly, ensuring that representation warping is strictly error-decreasing (dissipative) under linearization.
    \item \textbf{Feedback-Driven State Correction (RASC):} We introduce Representation-Agreement State Correction (RASC), which leverages representation-space coordinate tracking to dynamically correct and override biased feedforward routing confidence, completely neutralizing state-locking under persistent router corruption.
    \item \textbf{Rigorous Empirical Benchmarking \& Transparency:} We evaluate L-ARC against simpler weight-smoothing and feedback-decay heuristics under practical early-stage static centroids (Setting A), transient failures (Setting C), and confident systematic routing bias (Setting D) in the 14-layer Coordinate Sandbox. We report latency profiling, paired t-tests, and gating rates with complete scientific transparency---demonstrating that while ECG-Reset drives the main accuracy gains (+5.14\% under failures), our active controller and RASC mechanism provide essential state-space and representation-space semantic protection.
\end{itemize}
